\chapter{The transformer recipe}
\label{ch:transformer}

\newthought{This chapter} walks through a standard decoder-only transformer
block end to end, the way you would actually code it. The next chapter zooms
into the attention sub-layer; the chapter after that derives why that
sub-layer is the cost driver.

\section{The block, in one picture}

A decoder-only transformer is a stack of \(L\) identical blocks. Each block
takes a tensor of shape \([n, d]\) — \(n\) tokens, each represented by a
vector of \(\dmodel = d\) dimensions — and returns a tensor of the same
shape. Inside the block, two sub-layers do the work:

\begin{enumerate}
\item \term{Self-attention}, which lets every token mix information from
every other (earlier) token.
\item A \term{position-wise feed-forward network} (a small MLP applied
independently to each token), which gives the model nonlinear capacity
beyond what attention alone provides.
\end{enumerate}

Each sub-layer is wrapped in a \term{residual connection} and a
\term{normalisation} layer:
\begin{align}
  z &\;\leftarrow\; x + \text{Attn}(\text{LN}(x)) \\
  y &\;\leftarrow\; z + \text{MLP}(\text{LN}(z)).
  \label{eq:block}
\end{align}
That is a complete block.\sidenote{Two flavours: \emph{post-LN} (the
original 2017 paper) puts the LN \emph{after} the residual sum;
\emph{pre-LN} (now standard) puts it \emph{before}, as written here.
Pre-LN trains more stably at depth.} The whole model wraps token embeddings
+ positional information as input, runs \(L\) of these blocks, applies a
final LN, and projects to logits over the vocabulary.

\begin{figure}
\centering
\begin{tikzpicture}[
  node distance=4mm,
  block/.style={draw, rounded corners, minimum width=4.0cm, minimum height=6mm, align=center, font=\small},
  res/.style={draw, circle, inner sep=1pt, font=\small},
  every node/.style={font=\small},
]
  \node[block] (in)  {input \([n, d]\)};
  \node[block, below=of in] (ln1) {LayerNorm};
  \node[block, below=of ln1] (att) {Self-Attention};
  \node[res,   below=of att] (s1)  {\(+\)};
  \node[block, below=of s1]  (ln2) {LayerNorm};
  \node[block, below=of ln2] (mlp) {Feed-Forward};
  \node[res,   below=of mlp] (s2)  {\(+\)};
  \node[block, below=of s2]  (out) {output \([n, d]\)};
  \draw[-{Latex[length=2mm]}] (in) -- (ln1);
  \draw[-{Latex[length=2mm]}] (ln1) -- (att);
  \draw[-{Latex[length=2mm]}] (att) -- (s1);
  \draw[-{Latex[length=2mm]}] (s1) -- (ln2);
  \draw[-{Latex[length=2mm]}] (ln2) -- (mlp);
  \draw[-{Latex[length=2mm]}] (mlp) -- (s2);
  \draw[-{Latex[length=2mm]}] (s2) -- (out);
  \draw[-{Latex[length=2mm]}, dashed] (in.east) -- ++(1.0,0) |- (s1.east);
  \draw[-{Latex[length=2mm]}, dashed] (s1.east) -- ++(1.0,0) |- (s2.east);
\end{tikzpicture}
\caption{One transformer block (pre-LN). Dashed lines are residual
connections; the two solid sub-blocks (attention and feed-forward) are where
all the actual work happens.}
\label{fig:block}
\end{figure}

\section{The pieces, named}

\paragraph{Embeddings.} Token ids \(\to\) vectors of dimension \(d\) via the
table from Chapter~\ref{ch:tokens}.

\paragraph{Positional information.} Attention is permutation-invariant — if
you shuffle the tokens, the output is the same. Real language is not. The
fix is to add positional information to the embeddings.\sidenote{Three
families have dominated. (i) Learned absolute positions (BERT, GPT-2):
add a learned vector per index. (ii) Sinusoidal (Vaswani 2017): use
hand-crafted sine waves at geometric frequencies. (iii) RoPE (rotary
position embeddings, \citealp{su2021rope}): apply a position-dependent
rotation inside the attention's Q and K, which is the de-facto modern
choice (Llama, Mistral, Gemma). RoPE composes cleanly with extrapolation
tricks like YaRN \citep{peng2023yarn}.}

\paragraph{Self-attention.} The hero of Chapter~\ref{ch:self-attention}.
Lets every token's representation be a weighted average of every other
token's value, where the weights are computed from the dot product of the
current token's \emph{query} with the others' \emph{keys}.

\paragraph{Feed-forward (MLP).} A two-layer MLP applied independently at
each position:
\begin{equation}
  \text{MLP}(x) \;=\; W_2\,\sigma(W_1 x + b_1) + b_2,
\end{equation}
with hidden dimension \(d_{\text{ff}}\), typically \(4d\) (so 16,384 for
\(d = 4{,}096\)). \(\sigma\) is GeLU, SwiGLU, or similar. This is where
most of the parameters live — the FFN is roughly \(8 d^2\) parameters per
layer, vs. \(4 d^2\) for attention's projection matrices.\sidenote{For a
7B model with \(d=4096\), \(L=32\): attention \(\approx 32 \cdot 4 d^2
\approx 2.1\)B params, FFN \(\approx 32 \cdot 8 d^2 \approx 4.3\)B params,
plus embeddings and norms. The FFN dominates the parameter count even
though attention dominates the asymptotic compute.}

\paragraph{LayerNorm / RMSNorm.} A normalisation that rescales each token's
hidden vector to have zero mean and unit variance (LN) or just unit norm
(RMSNorm). Standard since 2018, plus a small learned scale and shift.

\paragraph{Output projection.} A linear layer of shape \([d, V]\) followed
by a softmax (equation~\ref{eq:softmax}) to produce next-token probabilities.

\section{Why the transformer beat the LSTM}

Three reasons matter.

\begin{enumerate}
\item \textbf{Parallelism on the time axis.} Equation~\eqref{eq:rnn} forced
\(h_t\) to be computed before \(h_{t+1}\). Self-attention computes all
positions at once via a single matrix multiply; the GPU stays busy.
\item \textbf{No vanishing gradient over distance.} An LSTM's signal from
token 1 to token 1000 has to survive 999 multiplicative updates. In a
transformer there is a direct dot product between any two positions in any
layer.
\item \textbf{Easy depth scaling.} Residuals + LayerNorm let you stack 100+
identical blocks with no special handling. RNNs above ${\sim}4$ layers were
notoriously hard to train.
\end{enumerate}

\section{What it cost}

The transformer pays for those three wins with one bill: every layer's
attention sub-block computes an \(n \times n\) matrix of pairwise
similarities between tokens. Compute scales as \(n^2\). Memory for the
attention matrix scales as \(n^2\). The KV cache used at generation time
scales as \(n\) per layer, but with multiplicatively large constants.

For \(n\) of a few thousand — the regime of GPT-2 and GPT-3 — this was fine.
For \(n\) of one million — what 2026's agentic workloads want — it is the
defining problem of the field. Chapter~\ref{ch:why-quadratic} derives the
\(n^2\) cost carefully; Part~IV documents the engineering response;
Part~V covers the architectures that try to escape it.
